\documentclass[a4paper,10pt]{article}
\usepackage[utf8]{inputenc}
\usepackage{graphicx}
\usepackage{subfigure}
\usepackage{fancyhdr}
\usepackage{listings}
\usepackage{ gensymb }


%opening

\pagestyle{fancy} %eigener Seitenstil
\fancyhf{} %alle Kopf- und Fußzeilenfelder bereinigen
\fancyhead[L]{SEP Abgabegespräch SS2015/16} %Kopfzeile links
\fancyhead[C]{} %zentrierte Kopfzeile
\fancyhead[R]{\includegraphics[width=1cm]{tugraz_bw}} %Kopfzeile rechts
\renewcommand{\headrulewidth}{0.4pt} %obere Trennlinie
\fancyfoot[C]{} %Seitennummer
\renewcommand{\footrulewidth}{0.4pt} %untere Trennlinie


\lstset{language=C++, basicstyle=\ttfamily\tiny}  

\begin{document}

\section*{Simple File System}

Entwickle ein Programm welches ein Filesystem in einer anschaulichen form darstellt. Als input dient ein File welches die einzelnen Eintraege beschreibt. Wo bei ein Eintrag folgende Form hat:
\texttt{<Type> <Name> <Parent>}. Wobei der \texttt{Type} nur 'O' oder 'F' sein darf (Ordner oder File). Der Ordner 'root' soll immer exsitieren, dieser ist vorgegeben und hat als einziger keinen Parent. 

Eine Beispiel Eingabedatei \texttt{files.sys}:
\begin{lstlisting}[frame=single]
O noten root
F passwords secure
O utils root
O storage utils
F ls storage
F 3 noten
O secure utils
\end{lstlisting}

Das Programm soll diese Datei einlesen und diese in folgender form ausgeben:
\subsubsection*{Ausgabe}
\begin{verbatim}
root
    noten
        3
    utils
        storage
            ls
        secure
            passwords
\end{verbatim}

Die Reihenfolge soll hier keine Rolle spielen. Wichtig ist die Reihenfolge der Ordner wie sie zum Speichern hinzugefuegt werden.

Sollte ein Fehler (File not found, File/Ordner nicht in Struktur) auftreten soll
\begin{verbatim}
ERR
\end{verbatim}
ausgegeben werden.

\subsection*{Aufruf}
\texttt{./ag files.sys}

\subsection*{Bewertung}
\begin{itemize}
 \item Dokumentation und Programmierstil: 1 Punkt
 \item Einhalten der OOP Konzepte: 1 Punkt
 \item Struktur \& Korrektheit: 1 Punkt
\end{itemize}

\newpage


\end{document}
